\documentclass[10pt,letterpaper]{article}
\usepackage[letterpaper,margin=0.75in]{geometry}
\usepackage[T1]{fontenc}
\usepackage{lmodern}
\usepackage{microtype}
\usepackage{parskip}
\usepackage{amsmath}
\usepackage{amssymb}
\usepackage{booktabs}
\usepackage{tabularx}
\usepackage{enumitem}
\usepackage[hidelinks]{hyperref}
\usepackage{xurl}
\setlist{nosep,leftmargin=1.4em}
\newlist{todo}{itemize}{2}
\setlist[todo]{label=$\square$,nosep,leftmargin=1.6em}
\usepackage{titlesec}
\titlespacing*{\section}{0pt}{10pt}{3pt}
\titlespacing*{\subsection}{0pt}{8pt}{2pt}
\titleformat*{\section}{\large\bfseries}
\titleformat*{\subsection}{\normalsize\bfseries}

\begin{document}

{\LARGE\bfseries BATS: Stage 1 Implementation Notes\par}
\vspace{2pt}
{\large Single-GPU tensor-core beamformer \hfill UWHPC \textperiodcentered{} Working draft, \today\par}
\vspace{-4pt}
\noindent\rule{\linewidth}{0.6pt}

Working notes for researching how to build Milestone~1 of the technical specification. They cover what stage~1 includes, what is already decided, and what still needs research.

\section*{Goal of Stage 1}

Milestone~1 is complete when:
\begin{enumerate}
  \item the single-GPU beamformer, the workload generator, and the reference implementations exist;
  \item the beamformer passes the correctness tests (spec, Numerical Correctness) within documented tolerances;
  \item its performance has been measured against the references.
\end{enumerate}

For each of the $F$ independent frequency channels $f$, the computation is
\[
  Y_f = W_f X_f, \qquad
  W_f \in \mathbb{C}^{B \times S}\ \text{(weights)}, \quad
  X_f \in \mathbb{C}^{S \times T}\ \text{(samples)}, \quad
  Y_f \in \mathbb{C}^{B \times T}\ \text{(beams)}.
\]
This is a \textbf{batched complex GEMM} with batch count $F$. The work is $8BSTF$ FLOPs and the arithmetic intensity is $8BST / \big(e\,(ST + BS + BT)\big)$, where $e$ is the number of bytes per complex value.

\section*{Components}

\begin{tabularx}{\linewidth}{@{}l l X@{}}
  \toprule
  \textbf{Component} & \textbf{Purpose} & \textbf{Notes} \\
  \midrule
  FP32 reference (cuBLAS) & Correctness and performance & Complex strided batch in true FP32 (\texttt{CUBLAS\_COMPUTE\_32F}, \texttt{NVIDIA\_TF32\_OVERRIDE=0}). One library call in the test harness; no kernel to write. FP32 error ($\sim 10^{-7}$) is far below every BATS precision. \\
  cuBLAS TF32 & Performance comparison & Same call with \texttt{CUBLAS\_COMPUTE\_32F\_FAST\_TF32}, like for like with the BATS TF32 kernel. BATS itself never calls cuBLAS. \\
  Tensor-Core Beamformer & Existing tensor-core reference & Oostrum et al., IPDPS 2025. \\
  BATS kernel (CUTLASS 3.x) & What we are building & TF32 real GEMM on Hopper \texttt{wgmma}, using the block formulation for complex arithmetic, as a strided batch over channels. \\
  Workload generator & Shared synthetic inputs & Random data, plane-wave sources from known directions, edge cases. \\
  Correctness tests & Compare against reference & Relative Frobenius error plus maximum elementwise error. \\
  Benchmark harness & Repeatable timing & Records dimensions, hardware, software versions, and timings. \\
  \bottomrule
\end{tabularx}

\section*{Decisions So Far}

\begin{itemize}
  \item \textbf{Language and build:} CUDA/C++, C++23, CMake. clangd is configured in \texttt{.clangd}, formatting follows LLVM style (\texttt{.clang-format}), and \texttt{.clang-tidy} enables the bugprone, modernize, performance, and readability checks.
  \item \textbf{Hardware:} NVIDIA Hopper (\texttt{sm\_90a}: H100, H200) and Blackwell (\texttt{sm\_100a}: B200, GB200). Both need architecture-specific (\texttt{a}) builds to use their tensor-core instructions.
  \item \textbf{Hopper only to start.} Stage~1 targets Hopper only. Blackwell is added later as a second implementation; an \texttt{sm\_90a} binary does not run on Blackwell, so it is a new kernel, not a recompile.
  \item \textbf{CUTLASS 3.x only.} CUTLASS v4.8.1, using only its C++ 3.x API (CuTe, collective builders, \path{cutlass::gemm::kernel::GemmUniversal} with \path{GemmUniversalAdapter}). No CUTLASS 2.x classes, no hand-written PTX, and no cuBLAS inside BATS. The spec still says ``CUTLASS or direct tensor-core instructions'' (\texttt{approach.tex}); update it.
  \item \textbf{FP32 baseline, TF32 first kernel.} The baseline (correctness reference and speed reference) is FP32, computed by cuBLAS. The first BATS kernel uses plain TF32 (FP32 storage, inputs rounded to TF32, FP32 accumulation) on Hopper \texttt{wgmma}. No FP32 or 3$\times$TF32 emulation. TF32 rounds each input to about $5 \times 10^{-4}$ relative error (FP32: $6 \times 10^{-8}$); a rough estimate (assuming rounding errors are random and partly cancel) puts output error around $10^{-4}$ relative, roughly 10$\times$ below 8-bit and 100$\times$ below 4-bit sensor quantization. Related work supports this: the Tensor-Core Beamformer uses 16-bit and 1-bit precision, and the Tensor-Core Correlator shows low precision is sufficient for closely related radio-astronomy computations; FP16 has the same 10-bit mantissa as TF32. TF32 is therefore \emph{expected} to be accurate enough, but this is not yet verified: the FP32 reference must measure the actual error, and ``enough'' depends on application requirements (dynamic range, sidelobe level) that still have to be set. After TF32 works, try lower precisions through the same alias and kernel: FP16 and BF16, then FP8, all on Hopper. FP4 (and FP6) need Blackwell and wait until it is added. FP64 is not used anywhere: the FP32 cuBLAS reference is accurate enough for every precision tested.
  \item \textbf{Complex arithmetic by the block formulation.} CUTLASS 3.x has no complex kernel for Hopper, so each channel's complex product is computed as one real GEMM:
  \[
    \begin{bmatrix} Y_r \\ Y_i \end{bmatrix}
    =
    \begin{bmatrix} W_r & -W_i \\ W_i & W_r \end{bmatrix}
    \begin{bmatrix} X_r \\ X_i \end{bmatrix},
    \qquad
    (2B \times T) = (2B \times 2S)(2S \times T).
  \]
  If each channel stores its real plane followed by its imaginary plane, $[X_r; X_i]$ is already one $2S \times T$ matrix and $[Y_r; Y_i]$ comes out planar, so inputs and outputs need no conversion. Only the $2B \times 2S$ block weight matrix is built, once, and reused across sample blocks (counted as planning cost). The work equals the 4M method.
  \item \textbf{No transpose.} The kernel computes $Y_f = W_f X_f$ as written, with $M = 2B$, $N = T$, $K = 2S$ in the block formulation. The transposed form $Y^{\mathsf T} = X^{\mathsf T} W^{\mathsf T}$ is not used. Known cost: the MMA's $M$ dimension is at least 64, so with fewer than 32 beams ($2B < 64$) tiles are partly empty (Workload~A).
  \item \textbf{Precision through a type alias.} The working precision is chosen in one place with a type alias (\texttt{using Real = float;} for TF32) so precisions can be switched and compared quickly. The alias selects a \textbf{precision traits} struct that bundles the storage type, MMA input type, accumulator type, and kernel configuration. The reference, generator, and tests are written generically over the alias; only the kernel choice is precision specific. Later precisions (FP16, BF16, FP8) reuse the same block-formulation kernel with different element types.
  \item \textbf{Strided batching through zero-cost views.} The $F$ channels are stored back to back, so channel $f$ starts at \texttt{base + f $\times$ stride}. CUTLASS 3.x takes the batch as the $L$ mode of the problem shape $(M, N, K, L)$ with a stride for each operand, so one launch covers all channels. All code indexes the data through \texttt{cuda::std::mdspan} views (CCCL, usable on host and device), e.g.\ \texttt{x(f, i, t)}, which compile to the same pointer arithmetic as hand-written indexing. CUTLASS strides come straight from the views.
  \item \textbf{Architecture behind an interface.} The beamformer is called through an architecture-independent interface, with the kernel chosen from the device's compute capability at run time. Adding Blackwell then means adding a second implementation, and the distributed planner can later treat each GPU's architecture as part of the topology it reads.
  \item \textbf{Input layout: $[f][t][s]$, split per sample.} Within each channel, every time sample stores the real parts of all $S$ sensors followed by their imaginary parts. This is exactly $[X_r; X_i]$ as a $2S \times T$ column-major matrix with the sensor ($K$) dimension contiguous, so TF32 \texttt{wgmma} reads it with no transpose or copy. Station data usually arrives time-ordered per sensor, so ingest needs one reorder, counted as a layout cost. cuBLAS needs interleaved complex, so the test harness converts for the reference (a test-only cost).
  \item \textbf{Input data:} weights and samples are supplied. Channelization, calibration, and weight generation are out of scope.
  \item \textbf{CUTLASS as a git submodule,} pinned to tag \texttt{v4.8.1}.
  \item \textbf{GoogleTest} for correctness tests, integrated with CTest.
  \item \textbf{Layout costs count:} any packing, transposition, padding, or copies a method needs are included in its timing (spec, Evaluation).
\end{itemize}

\section*{Hopper Versus Blackwell Kernels}

Both architectures load tiles with TMA, synchronize with \texttt{mbarrier}, group CTAs into clusters, and use warp-specialized persistent kernels. The tensor-core instruction and where its accumulator lives are different, so each needs its own CUTLASS kernel instantiation.

\begin{tabularx}{\linewidth}{@{}l X X@{}}
  \toprule
  & \textbf{Hopper (\texttt{sm\_90a})} & \textbf{Blackwell (\texttt{sm\_100a})} \\
  \midrule
  MMA instruction & \texttt{wgmma}, issued by a warpgroup (4 warps) & \texttt{tcgen05.mma}, issued by a single thread \\
  Operands & A from registers or shared memory; B from shared memory & A and B from shared memory; A may also come from tensor memory \\
  Accumulator & Registers of the issuing warpgroup & Tensor memory (TMEM), a dedicated 256\,KB per-SM store, read out with \texttt{tcgen05.ld} \\
  Tile shape & $M = 64$ per warpgroup; $N$ from 8 to 256 in steps of 8 & $M = 64$ or 128 on one SM, 128 or 256 across a 2-SM CTA pair; $N$ up to 256 \\
  Multi-SM MMA & No; clusters share loads through TMA multicast only & Yes; one MMA can span two SMs (\texttt{cta\_group::2}) \\
  Epilogue & Uses the same registers as the accumulator, so it competes with the main loop & Accumulator stays in TMEM, so the next tile's MMA can overlap the current tile's epilogue \\
  Scheduling & Persistent kernels with a static tile scheduler & Adds cluster launch control for dynamic tile scheduling \\
  Precisions (\texttt{wgmma}/\texttt{tcgen05}) & FP16, BF16, TF32, FP8, INT8; \textbf{no FP64} & Adds FP6, FP4, and block-scaled MXFP8/6/4 and NVFP4 \\
  CUTLASS & Arch tag \texttt{Sm90}; ping-pong and cooperative warp-specialized schedules & Arch tag \texttt{Sm100}; 1-SM and 2-SM schedules \\
  \bottomrule
\end{tabularx}

\textbf{Consequences for BATS} (to verify during research):
\begin{itemize}
  \item \textbf{Small $B$.} The MMA's $M$ dimension has a minimum of 64. Since we do not transpose, $M = 2B$, so fewer than 32 beams leaves tiles partly empty. Measure how much this costs Workload~A.
  \item \textbf{Shared code.} Data layout, workload generation, tests, and benchmarking are architecture independent. Only the kernel instantiations (mainloop and epilogue schedules, tile and cluster shapes) differ.
\end{itemize}

\section*{CUTLASS Pinning and Hopper Support}

Researched against the CUTLASS repository on 14 November 2026.

\begin{itemize}
  \item \textbf{Release:} pin tag \texttt{v4.8.1} (15 October 2026), the newest release. Its C++ code generation is tuned for CUDA~13.4. CUTLASS itself needs C++17; \texttt{sm\_90a} needs CUDA~12.0 or later, and adding Blackwell later needs 12.8 or later. Use CUDA~13.4 unless the cluster has something older, and record the exact toolkit and driver versions with every benchmark.
  \item \textbf{Known issues in 4.8.1:} a toolkit code-generation bug affects some packed FP8 and mixed-input instructions, mainly on Rubin (\texttt{sm\_107a}). The TF32 path should not be affected; recheck before trying FP8.
  \item \textbf{Header-only:} add \texttt{include/} (and \texttt{tools/util/include/} for helpers) to the include path. No CUTLASS build step is needed. Building CUTLASS's own tests and examples needs \texttt{-DCUTLASS\_NVCC\_ARCHS=90a}.
  \item \textbf{TF32 on Hopper:} CuTe provides TF32 \texttt{wgmma} atoms (e.g.\ \texttt{MMA\_64x256x8\_F32TF32TF32\_SS\_TN}). The SM90 collective builder maps \texttt{float} inputs to TF32 MMA automatically.
  \item \textbf{Layouts:} the TF32 \texttt{wgmma} instruction is natively ``TN'' (both operands K-major, i.e.\ the reduction dimension $S$ contiguous). The SM90 builder accepts other layouts by swapping A and B or by transposing B in shared memory (\path{is_swapAB}, \path{is_warpspecialized_transpose_B} in \path{sm90_gmma_builder.inl}), so any layout works, but K-major is expected to be fastest.
  \item \textbf{FP64 on Hopper:} CuTe has FP64 and complex-FP64 \texttt{mma.sync} atoms (\path{SM90_16x8x4_C64C64C64C64_TN} and 16x8x8 / 16x8x16 variants), but no FP64 \texttt{wgmma} and no tested 3.x Hopper FP64 GEMM. Not used.
  \item \textbf{Complex on Hopper:} no 3.x complex kernel for any reduced precision; hence the block formulation. (Blackwell does have a native 3.x planar complex GEMM, \path{sm100_mma_array_warpspecialized_planar_complex}, for FP16/BF16; consider it when Blackwell is added.)
\end{itemize}

\textbf{Starting points:}
\begin{itemize}
  \item CUTLASS's Hopper GEMM examples (\path{examples/48_hopper_warp_specialized_gemm} and nearby) and the SM90 TF32 device tests (\path{test/unit/gemm/device/sm90_gemm_tf32_*}) for a 3.x \texttt{CollectiveBuilder} with arch tag \texttt{Sm90} and \texttt{float} inputs.
  \item The CuTe tutorials under \path{examples/cute/tutorial} for building layouts and tensors.
\end{itemize}

\section*{Open Research Questions}

\subsection*{CUTLASS}
\begin{todo}
  \item \textbf{Cluster toolkit.} Which CUDA toolkit and driver the cluster has; whether CUDA~13.4 is available.
  \item \textbf{Batched problem shape.} Set up the $(M, N, K, L)$ problem shape with $L = F$ and per-operand strides (\texttt{cute::Stride<\ldots, int64\_t>}) from the \texttt{mdspan} views. Confirm that the batch stride may be any value, including the padded per-channel size.
  \item \textbf{Operand layout.} With $Y = W X$, the block weight matrix stored row-major and the $[f][t][s]$ samples are both K-major. Confirm the builder selects the native TN path with no swap or shared-memory transpose.
  \item \textbf{Block weight matrix.} Where to build $\left[\begin{smallmatrix} W_r & -W_i \\ W_i & W_r \end{smallmatrix}\right]$ (host or a small GPU kernel), its memory cost ($4BS$ real values per channel, twice the complex weights), and how to count its cost as planning.
  \item \textbf{Kernel schedule.} Compare ping-pong and cooperative warp-specialized schedules, and the persistent tile scheduler, for Workload~A and Workload~B shapes.
  \item \textbf{Tile-shape selection.} Which tile and cluster shapes to instantiate, and how to handle small $B$ or $S$ without wasting partly empty tiles.
  \item \textbf{Dimensions that are not tile multiples.} How CUTLASS 3.x handles them: predication through TMA bounds, and alignment requirements (TMA needs 16-byte-aligned strides, i.e.\ multiples of 4 \texttt{float}s).
  \item \textbf{Integration.} Add CUTLASS as a git submodule at \texttt{third\_party/cutlass}; keep CUTLASS kernels in separate \texttt{.cu} files to limit rebuild times; set \texttt{CMAKE\_CUDA\_ARCHITECTURES} to \texttt{90a} (add \texttt{100a} with Blackwell support); check that clangd parses CUTLASS headers with the current \texttt{.clangd} flags.
  \item \textbf{C++23 with nvcc.} Confirm that nvcc, CUTLASS, and \texttt{cuda::std::mdspan} build cleanly with \texttt{-std=c++23}; fall back to C++20 for \texttt{.cu} files if needed.
\end{todo}

\subsection*{Precision}
\begin{todo}
  \item Design the precision traits: storage type, MMA input type, accumulator type, kernel configuration, and error tolerance, selected at compile time by the alias.
  \item Measure TF32 error against the FP32 reference for realistic input ranges, and check whether CUTLASS rounds or truncates FP32 to TF32 (truncation roughly doubles the error).
  \item Whether the 3M (Gauss) variant, built from three real GEMMs, is accurate enough to be worth its saving.
  \item Lower precisions on Hopper, in order: FP16 and BF16 (about 2$\times$ TF32 throughput), then FP8 E4M3/E5M2 (about 4$\times$). All accumulate in FP32. FP8 has only 2--3 mantissa bits, enough for 4-bit samples but coarse for weights; measure its effect on beam gain and sidelobes, and consider keeping weights at higher precision if CUTLASS's mixed-input kernels allow it.
  \item FP4 and FP6 (block-scaled MXFP8/6/4, NVFP4) are Blackwell only. Try them when Blackwell is added.
  \item Error tolerance for each precision, justified against typical sensor-data bit widths (often 4 to 8 bits per component).
\end{todo}

\subsection*{Baselines}
\begin{todo}
  \item \textbf{cuBLAS.} \texttt{cublasGemmStridedBatchedEx} on complex FP32 data: with \texttt{CUBLAS\_COMPUTE\_32F} as the correctness reference, and with \texttt{CUBLAS\_COMPUTE\_32F\_FAST\_TF32} as the like-for-like speed comparison. Verify the reference really runs in FP32 (math mode, \texttt{NVIDIA\_TF32\_OVERRIDE=0}). cuBLAS takes \textbf{interleaved} complex data, so the harness converts from the BATS layout; report that conversion separately.
  \item \textbf{Tensor-Core Beamformer.} How to build it, which input layouts and precisions it supports, and what conversion our generated data needs.
\end{todo}

\subsection*{Correctness Testing}
\begin{todo}
  \item Metrics: relative error $\lVert Y - Y_{\mathrm{ref}} \rVert_F / \lVert Y_{\mathrm{ref}} \rVert_F$ and maximum elementwise error.
  \item Check NaN, infinity, and integer overflow separately.
  \item Test cases:
  \begin{itemize}
    \item random data;
    \item plane-wave sources from known directions, where the steered beam should recover the source with the expected gain;
    \item dimensions that are not multiples of the tile sizes;
    \item $B = 1$ and $S = 1$.
  \end{itemize}
  \item GoogleTest with \texttt{gtest\_discover\_tests}; parameterized tests over dimensions and precisions.
\end{todo}

\subsection*{Benchmarking}
\begin{todo}
  \item Time with CUDA events, after warm-up runs, over repeated trials; report a stated statistic and a measure of variability.
  \item Report kernel-only time and end-to-end time across the stated input and output boundaries separately.
  \item Report throughput as beam samples per second and as achieved FLOP/s relative to the hardware TF32 peak.
  \item Use Nsight Compute to check tensor-core utilization and achieved memory bandwidth.
  \item Decide the stage~1 parameter sweep: ranges for $B$, $S$, $T$, and $F$ covering both Workload~A (few beams, memory bound) and Workload~B (many beams, compute bound).
\end{todo}

\subsection*{Workload Generator and Data Format}
\begin{todo}
  \item Define padding and alignment for the $[f][t][s]$ split-per-sample layout (TMA needs 16-byte-aligned strides, so $2S$ should be a multiple of 4).
  \item Define the \texttt{mdspan} view types (extents, \texttt{layout\_right} / \texttt{layout\_left} / \texttt{layout\_stride}) and use them in the generator, the reference, and all implementations.
  \item Plane-wave generator: array geometry, source directions, and steering weights $w_{f,b,i} = a_{b,i}\, e^{\,j 2\pi f \tau_{b,i}}$.
  \item Generate on the host or the GPU? Use seeded randomness so runs are reproducible.
\end{todo}

\section*{Suggested Order}
\begin{enumerate}
  \item CMake project skeleton with CUDA, GoogleTest, and CUTLASS v4.8.1 (git submodule) building for \texttt{sm\_90a}, and a CUTLASS 3.x Hopper TF32 GEMM example running on our GPU.
  \item Precision type alias and traits, \texttt{mdspan} view types for the strided batch, and workload generator.
  \item FP32 cuBLAS reference (with layout conversion) and GoogleTest correctness harness.
  \item Block weight-matrix construction.
  \item TF32 block-formulation kernel: CUTLASS 3.x real GEMM, batched over channels through the $L$ mode, fed from the views and passing the correctness tests.
  \item cuBLAS TF32 speed comparison.
  \item Benchmark harness and a first sweep.
  \item FP16 and BF16, then FP8, through the same alias and kernel, each checked against the FP32 reference.
  \item Tensor-Core Beamformer comparison.
  \item Tile-shape, schedule, and precision tuning.
  \item Blackwell (\texttt{sm\_100a}) kernel behind the same interface, passing the same tests.
\end{enumerate}

\section*{Starting References}
\begin{itemize}
  \item CUTLASS repository, documentation, and examples: \url{https://github.com/NVIDIA/cutlass}
  \item PTX ISA (\texttt{wgmma}, \texttt{tcgen05}): \url{https://docs.nvidia.com/cuda/parallel-thread-execution/}
  \item \texttt{cuda::std::mdspan} (CCCL): \url{https://nvidia.github.io/cccl/}
  \item cuBLAS documentation: \url{https://docs.nvidia.com/cuda/cublas/}
  \item Tensor-Core Beamformer: Oostrum et al., IPDPS 2025 (\texttt{oostrum2025tcbf} in \texttt{references.bib})
  \item Error bounds for the 3M method: Higham, 1992 (\texttt{higham1992complex} in \texttt{references.bib})
\end{itemize}

\end{document}
